%& -shell-escape
% !TeX root~=~thesis-abstract.tex
% !TeX encoding~=~UTF-8
% !TeX program~=~XeLaTeX

\documentclass{.local/lib/classes/ndvr-super-article-standalone}

\begin{document}

    Рассмотрим кусочно-агрегатную аппроксимацию 
    в~алгоритме $\iSAX$.
    Пусть дан исходный временной ряд $C~=~(c_{i})^{n}_{i=1}$.
    \begin{enumerate}
    \item 
        Исходный ряд $C$ разбивается на $w$ отрезков.
        На его основе строится 
        новый ряд $\overline{C}~=~(\overline{c}_{i})^{w}_{i=1}$.
        Каждый элемент ряда вычисляется по следующей формуле:
        \[
            \overline{c}_{i} 
               ~=~\dfrac{w}{n} \cdot
                    \sum
                        \limits%
                            ^{\frac{n}{w} \cdot i}% 
                            _{j~=~\frac{n}{w}\cdot (i-1) + 1}%
                    c_{j}
        \]  
        % Для исходного ряда 
        % $C~=~\left\{2, 4, 7, 9, 5, 3, 4, 6\right\}$:
        % \[
        %     \begin{array}{rcccl}
        %         \overline{C} 
        %         &~=~& (\overline{c}_{i})^{4}_{i=1}    
        %         &~=~& \left\{3, 8, 4, 5\right\}; \\
        %     \end{array}
        % \]
    \item 
        Значения ряда $\overline{C}$ подвергаются дискретизации 
        в заданный набор значений мощности $s$.
        На основе дискретизированных значений строится ряд  
        $D^{s}~=~(d^{s}_{i})^{w}_{i=1}$.
        % Например:
        % \[
        %     \begin{array}{rcccl}
        %         C 
        %         &~=~& (c_{i})^{8}_{i=1}               
        %         &~=~& \left\{2, 4, 7, 9, 5, 3, 4, 6\right\}; \\
        %         \overline{C} 
        %         &~=~& (\overline{c}_{i})^{4}_{i=1}    
        %         &~=~& \left\{3, 8, 4, 5\right\}; \\
        %         D^{4} 
        %         &~=~& (d^{4}_{i})^{4}_{i=1}           
        %         &~=~& \left\{
        %             \mathrm{a}, 
        %             \mathrm{г}, 
        %             \mathrm{б}, 
        %             \mathrm{в}
        %         \right\}; \\
        %     \end{array}
        % \]
        % В оригинальном алгоритме в качестве элементов ряда
        % используются буквы латинского алфавита.
        % Для наглядности в примере использованы символы 
        % кириллицы. 
    \item 
        Каждый элемент ряда $D$ переводится 
        в свою бинарную форму.
        Из полученных форм строится ряд
        $B^{s}~=~(b^{s}_{i})^{w}_{i=1}$, 
        такой ряд обозначают\\
        $\SAX(C, w, s)~=~B^{s}$.
        % Например:
        % \[
        %     \begin{array}{rcccl}
        %         C 
        %         &~=~& (c_{i})^{8}_{i=1}               
        %         &~=~& \left\{2, 4, 7, 9, 5, 3, 4, 6\right\}; \\
        %         \overline{C} 
        %         &~=~& (\overline{c}_{i})^{4}_{i=1}    
        %         &~=~& \left\{3, 8, 4, 5\right\}; \\
        %         D^{4} 
        %         &~=~& (d^{4}_{i})^{4}_{i=1}           
        %         &~=~& \left\{
        %             \mathrm{a}, 
        %             \mathrm{г}, 
        %             \mathrm{б}, 
        %             \mathrm{в}
        %         \right\}; \\
        %         \SAX(C, 4, 4) 
        %         &~=~& (b^{2}_{i})^{4}_{i=1}           
        %         &~=~
        %         & \left\{
        %             {(0,0)}, 
        %             {(1,1)}, 
        %             {(0,1)}, 
        %             {(1,0)} 
        %         \right\}; \\
        %         \SAX
        %         (C, 4, 2) 
        %         &~=~& (b^{2}_{i})^{4}_{i=1}           
        %         &~=~
        %         & \left\{
        %             {(0)}, 
        %             {(1)}, 
        %             {(0)}, 
        %             {(1)}
        %         \right\}; \\
        %         \SAX(C, 4, 8) 
        %         &~=~& (b^{8}_{i})^{4}_{i=1}  
        %         &~=~         
        %         & \left\{
        %             {(0,0,0)}, 
        %             {(1,0,1)}, 
        %             {(0,1,0)}, 
        %             {(1,0,0)}
        %         \right\}; \\
        %     \end{array}
        % \]
    \item 
        Каждый элемент ряда $\SAX(C, w, s)$
        обладает длиной $l~=~\log_{2}{s}$. 
        При больших $w$ такое представление избыточно. 
        Кроме того, сравнивать временные ряды 
        с разным $s$ не удобно.
        Для решения этих проблем на основе
        ряда $(b^{l}_{i,q})^{w}_{i=1}$
        строится ряд \[
            \iSAX(C, w, l) 
               ~=~(s^{l_i}_{i})^{w}_{i=1}                     
                   ~=~(b^{l_i}_{i}, l_i)^{w}_{i=1}|_{l_i \le l}
        \] 
        % Например:
        % \[
        %     \begin{array}{rcccl}
        %         C 
        %         &~=~& (c_{i})^{8}_{i=1}               
        %         &~=~& \left\{2, 4, 7, 9, 5, 3, 4, 6\right\}; \\
        %         \overline{C} 
        %         &~=~& (\overline{c}_{i})^{4}_{i=1}    
        %         &~=~& \left\{3, 8, 4, 5\right\}; \\
        %         D^{4} 
        %         &~=~& (d^{4}_{i})^{4}_{i=1}           
        %         &~=~& \left\{
        %             \mathrm{a}, 
        %             \mathrm{г}, 
        %             \mathrm{б}, 
        %             \mathrm{в}
        %         \right\}; \\
        %         \SAX(C, 4, 4) 
        %         &~=~& (b^{2}_{i})^{4}_{i=1}           
        %         &~=~
        %         & \left\{
        %             {(0,0)}, 
        %             {(1,1)}, 
        %             {(0,1)}, 
        %             {(1,0)} 
        %         \right\}; \\
        %         \iSAX(C, 4, 4) 
        %         &~=~& (s_{i})^{4}_{i=1}       
        %         &~=~       
        %         & \left\{
        %             {(0)}^{2}, 
        %             {(1,1)}^{4}, 
        %             {(1)}^{2}, 
        %             {(1,0)}^{4}
        %         \right\}; \\
        %     \end{array}
        % \]
    \end{enumerate}
    Пример для исходного ряда 
    $\left\{2, 4, 7, 9, 5, 3, 4, 6\right\}$
    имеет вид:
    \[
        \begin{array}{rcccl}
            C 
            &~=~& (c_{i})^{8}_{i=1}               
            &~=~& \left\{2, 4, 7, 9, 5, 3, 4, 6\right\}; \\
            \overline{C} 
            &~=~& (\overline{c}_{i})^{4}_{i=1}    
            &~=~& \left\{3, 8, 4, 5\right\}; \\
            D^{4} 
            &~=~& (d^{4}_{i})^{4}_{i=1}           
            &~=~& \left\{
                \mathrm{a}, 
                \mathrm{г}, 
                \mathrm{б}, 
                \mathrm{в}
            \right\}; \\
            \SAX(C, 4, 4) 
            &~=~& (b^{2}_{i})^{4}_{i=1}           
            &~=~
            & \left\{
                {(0,0)}, 
                {(1,1)}, 
                {(0,1)}, 
                {(1,0)} 
            \right\}; \\
            \iSAX(C, 4, 4) 
            &~=~& (s_{i})^{4}_{i=1}       
            &~=~       
            & \left\{
                {(0)}^{2}, 
                {(1,1)}^{4}, 
                {(1)}^{2}, 
                {(1,0)}^{4}
            \right\}; \\
        \end{array}
    \]
        
\end{document}
